% !TeX root = ../papers/nc-p3.tex
\section{十五次元成分の三次曲面対と残余八点}
\label{sec:marked-cubic-surface-pairs}
\begin{theorem}\label{thm:marked-cubic-surface-pairs}
十五次元成分の稠密開集合で、二十点のいずれを印にしても以下が成立する。
位置 $0,2$ を反対対とし、印を入れた二つの三文字テンソルから、同じ
$\mathbb P^2\times\mathbb P^2$ 内の順序付き二曲面を得る。いずれも滑らかな三次曲面で、
二つの平面への射影を通じて五次の平面Cremona変換のグラフになる。
両曲面は長さ27で交わる。そのうち19点が元の印を除く点列の射影で、各々交点重複度は1である。
残余は長さ8である。二曲面の順序と二つの平面の順序を記録した射影同値類は、
その印を持つ窓の添字保持同型の完全不変量である。もう一つの反対対にも同じ結論が通る。
\end{theorem}

\begin{proof}[二曲面とその逆構成]
定理\ref{thm:marked-point-commutators}の四組のうち $G,F$ を使う。$T=\langle I,G_1,G_2,G_3\rangle$,
$S=\langle I,F_1,F_2,F_3\rangle$ は同じ三行三列の行列空間に入る二つの四平面である。
各 $E=T,S$ について $\det Z=0$ を $\mathbb P E$ 内の三次曲面とし、
$Z$ の左核 $l$ と右核 $r$ を二平面の積に写す。用いる開条件とその非空性は下で検証する。
滑らかな三次曲面には階数1の点がない。$l,r$ をともに核に持つ行列の空間は
高々一次元である。二次元なら、その鉛筆の二行二列の行列式に射影的零点があり階数1以下になるためである。
従って像上で行列を線型方程式の唯一の核として復元でき、像への写像は同型である。

$l$ を固定すると $lZ=0$ は $E$ に三式であり常に解を持つ。
一般ファイバーは零次元なので一点である。右核も同様で、二つの平面射影は双有理である。
両引戻し直線束を $L,R$、三次曲面の超平面類を $H$ と書くと $L^2=R^2=1$。
余因子行列の二次式に基点はなく、左右核のSegre写像を与えるので $L+R=2H$。
$H^2=3$ により $LR=5$、したがって平面間の双有理写像の次数は5である。

逆にこの埋め込まれた曲面から $E$ を元の行列空間の部分空間として取り戻せる。
反標準束は $H$ であるから、それによる埋込みは元の $\mathbb P^3$ 表示を座標変更まで与える。
各点で指定された左右核を持つ一次形式行列をこの四つの射影座標で解く。
元の $M$ に加え $M'$ も解なら、$M'\operatorname{adj}M$ の各成分が曲面上で零になる。
各成分は三次式なので $M'\operatorname{adj}M=C\det M$ で $C$ は定数行列、従って $M'=CM$。
左核の一致条件は $l C$ が $l$ の倍であることを全ての $l$ について与える。従って $C$ はスカラーである。
一次形式の係数の張る四平面は反標準座標の選び方にも依存しない。

二曲面から復元した $T,S$ は可逆な一行列の直線で交わる。これを単位行列として二つの三次元空間を双対に同定すれば、
残る基底変更は共役と倍率である。両四平面の跡零部分を $T_0,S_0$ とする。
混合交換子 $T_0\otimes S_0\to\mathfrak{sl}_3$ の階数は8、核線はもとの表示で
$\sum G_i\otimes F_i$ である。これは定理\ref{thm:marked-point-commutators}の式の跡であり、
非退化なので両三平面の双対基底を共通倍率まで指定する。あとは得られた
$[G_a,F_c]$ に定理\ref{thm:marked-mixed-bracket-reconstruction}の逆写像を用いればよい。
両側単射の開条件は全二十印で成立する（定理\ref{thm:marked-twenty-rational-invariants}）。
この復元の基底変更と倍率は同定理の印付き同型作用に入る。逆方向は構成の自然性から従う。
\end{proof}

\begin{proof}[交点の長さと開条件の算術的検証]
第一平面の点 $l$ を変数とし、$E$ の基底 $M_0,\ldots,M_3$ の四行 $l M_i$ の符号付き三次小行列式を $a_i$ とする。
$Z_E(l)=\sum a_i M_i$ の右核が五次写像 $f_E(l)$ である。実際、$Z_E$ の二行の外積は
除いた行の $l$ 座標倍の同一の五次ベクトルとなる。これは $lZ_E=0$ と座標環の一意分解から従う。
$a_i$ の全てが零のときだけ第一射影のファイバーが正次元となる。

二曲面の第一射影の正次元ファイバーの基点は互いに素である。
さらに第一平面の一般直線上で二つの五次写像の像が一致する点はない。
従って曲面の交叉には第一平面の曲線を覆う曲線も、第一平面の一点上の曲線もない。
各曲面の二つの射影次数と混合次数は $1,5,1$ であるから、積のChow環での交点数は
$1+25+1=27$ である。滑らかな四次元多様体の中の滑らかな余次元2の二曲面の固有交叉なので、
これはスキームの長さである。

他の点列の位置1,3の座標を印から射影すれば、その二点は
$w(p_0,-,x_2,-)$ と $w(x_0,-,p_2,-)$ の双方の左右核である。
印を除く19個の第一射影は相異なる。一般位置での三点共線の不在は
系\ref{cor:klein-four-configurations}の標本とエタール特殊化から従う。
これらの点は両写像の定義域にあり、二つの五次ベクトルの外積の微分は各点で階数2を持つ。
従ってそれぞれ単純交点で、残余の長さは8である。

\path{computations/verify_marked_cubic_surface_pairs.py} は
定理\ref{thm:marked-twenty-rational-invariants}の実在する1009進持上げを14桁で用いる。
全二十印と四位置で三次行列式の偏微分のイデアルが五次斉次成分を張ることを検査し、滑らかさを得る。
二方向の各曲面対で、両方の $a_i$ が生成するイデアルは第一平面の八次斉次成分を張る。
直線 $l=(1,s,2+3s)$ に沿う二つの五次ベクトルの外積から二つの十次斉次二元式を取り、
そのSylvester行列が全階数であることを検査する。無限遠もこの斉次終結式に含まれる。
全ての点を持ち上げて他の19点での両写像の非消滅と二つの接方向の独立性を確認する。
四平面の和の次元7および交換子の階数8も全対で検査する。いずれも単元消去および
残余行列の共通1009因子による除算後に十分な精度で非零な小行列を得る検査である。
従って標数零のこの持上げで開条件が同時に成立し、既約成分の一般点にも成立する。
\end{proof}

\begin{corollary}
この稠密開部分では、定理\ref{thm:recurrent-point-line-flags}の指定点を持つ二平面と、
同じく指定点を持つその後続を同時に取れない。
\end{corollary}
\begin{proof}
出発点列の印を $a$、後続平面の印を $b$ とする。
$U\subset H(a_i)$ なので $U^\perp$ の独立な第二元を $q$ とすると、
$w(a_{i-1},q,b_{i+1},-)=0$ が必要である。$a=b$ なら反対行列は階数3で、
その左核は $a_i$ のみである。異なる場合は $a$ の滑らかな三次曲面の点なので
反対行列の階数は2で、左核は $a_i,b_i$ の張る平面である。
従って後続の左支持を与える行列は非零倍率を除いて $w(a_{i-1},b_i,-,-)$ である。
$b_i$ の印からの射影では上の三次小行列式 $a_j$ が全消滅しないので、この縮約は階数3である。
二平面を後続にすることはできない。
\end{proof}
